% Motivation and learning outcomes.
% The outcomes are the outcomes of Day 5 in the syllabus.
\begin{frame}{Why this day matters}
  \begin{itemize}
    \item A microphone and a camera give much more data than an accelerometer:
          16 000 samples each second, and 27 648 values for one small image.
    \item The model cannot take this data directly. Features make an audio
          signal small, and a small image and a narrow network make a vision
          model small.
    \item A device that listens all the time makes many decisions each second.
          The post-processing decides how many of them are false activations.
    \item In the lab you put a keyword spotting model and an image classifier
          on the XIAOML Kit, and you measure their latency and their RAM.
  \end{itemize}
\end{frame}

\begin{frame}{Learning outcomes}
  After this day you can:
  \begin{itemize}
    \item Compute a spectrogram and MFCC features from an audio signal.
    \item Describe the keyword spotting pipeline and the cascade design.
    \item Explain why peak activation memory limits vision models on a
          microcontroller.
    \item Deploy a keyword spotting model and an image classifier on the
          XIAOML Kit.
    \item Tune the post-processing to reduce false activations.
  \end{itemize}
\end{frame}
